% ============================================================================
% RASD --- F1tenth PoliMI car
% Source: docs/scope.md (working source until v1.00), docs/hardware/*, hardware/bom/master_bom.xlsx
% ============================================================================
\documentclass[11pt,a4paper]{article}

\input{preamble}
\input{doc}
\docbrand{polimi}

\docsetup{Requirements Analysis and Specification Document (RASD)}{F1tenth PoliMI}{The car: base platform and advanced topic}{0.01}{F1tenth PoliMI team}

\begin{document}
\docfrontmatter{
0.01 & 2026-10-25 & F1tenth PoliMI team & Initial draft, folded from \code{docs/scope.md} (21 Oct 2026), the hardware comparisons and the master BOM; \code{scope.md} at commit \code{425ee04} \\
}

% ============================================================================
\section{Executive Summary}
% ============================================================================
\subsection{Overview}
We build a 1/10-scale autonomous car (F1TENTH) and its ROS~2 software stack. The car runs indoors, in a room, on a closed track drawn on the floor with tape. It maps the room with a LiDAR, finds its pose in that map, sees the tape with a camera and drives laps without leaving the lines. A person can always stop it or take it over.

The same stack runs on a simulator, so every function can be tried before the car. The professor will use the car and the stack for the summer school exercises (esercitazioni).

\subsection{Project Objectives}
\begin{itemize}
  \item \textbf{Base platform (required):} a complete car and stack that works end to end: drivers, simulator, mapping, localization, line detection, path following and safety, with the documentation another team needs to use it.
  \item \textbf{One advanced topic:} racing-line planning (option A) or hybrid localization (option B) \tbd[OI-01].
  \item \textbf{Key deliverable:} the 30L delivery, what our team delivers to get 30 e lode: the base platform plus one advanced topic.
\end{itemize}

\subsection{Scope}
\textbf{In scope:} the car's hardware as selected in the master BOM, the ROS~2 stack on the car and the stack's use of the simulator, the safety features, the documentation.

\textbf{Out of scope here:} the simulator's internals (the gym, the bridge and our \code{polimi\_sim} extensions). They get their own RASD and TDD. This document only requires that the simulator stands in for the car (\cref{sec:fr-sim}).

\textbf{Reference, not required:} the summer school exercises (\cref{sec:exercises}).

\begin{takeaway}
\textbf{Key point.} The LiDAR maps the room, the camera sees the track. Tape is flat on the floor and the LiDAR scans a plane above it, so nothing the LiDAR sees marks the track. Every requirement below follows from that.
\end{takeaway}

% ============================================================================
\section{Requirements Analysis}
% ============================================================================
\subsection{Business Requirements}
\subsubsection{Business Goals}
\begin{enumerate}
  \item \textbf{30 e lode:} deliver a state-of-the-art base platform and one advanced topic.
  \item \textbf{Summer school:} the professor can teach the exercises on our car and stack.
  \item \textbf{Hand-over:} another team or the professor can set up, run and extend the platform without us.
\end{enumerate}

\subsubsection{Stakeholder Analysis}
\begin{tabularx}{\textwidth}{@{}L{3.5cm} Y@{}}
\toprule
\hd{Stakeholder} & \hd{Interest} \\
\midrule
Our team & Grade (30L): a platform that works end to end and one advanced topic that measurably improves on it. \\
Professor & A car and stack to teach the summer school exercises on; documentation to run it without us. \\
Summer school students & Exercises that run in the simulator and on the car, from the simplest to the most advanced. \\
Future teams & A platform to extend: clear interfaces, setup and design decisions written down. \\
\bottomrule
\end{tabularx}

\subsubsection{Legend}
\begin{itemize}
  \item \textbf{Phase:} \code{V1} base platform; \code{V2} advanced topic; \code{Ref} reference only (summer school exercises, \tbd[OI-12]).
  \item \textbf{Src:} \code{S} \code{docs/scope.md}; \code{H} hardware comparisons in \code{docs/hardware/}; \code{B} master BOM; \code{D} a decided open item; \code{C} \code{docs/software/car.md}.
  \item \textbf{IDs:} \code{FR-<DOM>-nnn} functional, \code{NFR-<CAT>-nnn} non-functional, \code{C-nnn} constraint, \code{A-nnn} assumption.
\end{itemize}

\subsection{Constraints and Assumptions}
\subsubsection{Environment}
\begin{tabularx}{\textwidth}{K}
\toprule
\hd{ID} & \hd{Constraint} & \hd{Src} \\
\midrule
\endhead
C-001 & The car runs indoors, in a room. & S \\
C-002 & The track is a closed loop drawn on the floor with tape. It has no walls or barriers: nothing along the track rises above the floor. & S \\
C-003 & The only things standing up from the floor are the room itself: its walls, furniture and anything else in it. & S \\
C-004 & The LiDAR cannot see the track: it scans a horizontal plane above the tape and sees the room only. & S \\
C-005 & Leaving the track is not a collision: nothing physical stops the car when it crosses a line. & S \\
\bottomrule
\end{tabularx}

\subsubsection{Hardware}
\begin{tabularx}{\textwidth}{K}
\toprule
\hd{ID} & \hd{Constraint} & \hd{Src} \\
\midrule
\endhead
C-010 & Chassis: Traxxas Ford Fiesta ST Rally VXL (74276-4), 1/10 scale, brushless, 4WD. End of life, available through unofficial retailers. & B \\
C-011 & Compute: NVIDIA Jetson Orin Nano Super with a 500\,GB NVMe SSD. & B \\
C-012 & LiDAR: Orbbec Pulsar SL450: 270\textdegree, up to 40\,Hz at 0.2\textdegree, Ethernet. It takes the Orin Nano's only Ethernet port. & B, H \\
C-013 & Camera: Orbbec Gemini 335L: global-shutter colour 1280\x800 up to 60\,fps, 94\textdegree\x68\textdegree, depth aligned to colour, IMU, USB~3. & B, H \\
C-014 & Motor and servo through a VESC 6 MkVI; power through the RoboRacer power board; manual control through an 8BitDo Ultimate 2C wireless joystick. & B \\
C-015 & The track material (ducts) is not bought: the track is tape. & B, S \\
\bottomrule
\end{tabularx}

\subsubsection{Software}
\begin{tabularx}{\textwidth}{K}
\toprule
\hd{ID} & \hd{Constraint} & \hd{Src} \\
\midrule
\endhead
C-020 & ROS~2 Jazzy on every image, Mac and car. Humble ends on 31 May 2027, before or during the summer school; Jazzy is supported to 31 May 2029 (OI-13, decided 21 Oct 2026). & D \\
C-021 & The stack runs in Docker, the same on the simulator and on the car. & S \\
C-022 & Drivers come from our fork of \code{f1tenth\_system} (branch \code{jazzy-polimi}); the LiDAR and the camera are driven by Orbbec's \code{orbbec\_camera}. & S, D \\
C-023 & The simulator is upstream's JAX gym (\code{dev-jax}) and bridge (\code{dev-jazzy}) with our extensions in our fork of \code{f1tenth\_gym\_ros} (OI-13). & D \\
\bottomrule
\end{tabularx}

\subsubsection{Assumptions}
\begin{tabularx}{\textwidth}{K}
\toprule
\hd{ID} & \hd{Assumption} & \hd{Src} \\
\midrule
\endhead
A-001 & The room is not too large, and it is mapped with SLAM on site. & S \\
A-002 & Floor material and colour are not defined yet; ideally they vary too. \tbd & S \\
A-003 & Tape colour and width are not defined yet. \tbd & S \\
\bottomrule
\end{tabularx}

\subsection{Functional Requirements}
\subsubsection{User Roles}
\begin{tabularx}{\textwidth}{@{}L{3.5cm} Y Y@{}}
\toprule
\hd{Role} & \hd{Description} & \hd{Responsibilities} \\
\midrule
Operator & The person next to the car, holding the joystick & Starts and stops runs, holds the deadman, takes over, presses the emergency stop \\
Developer & Team member or student writing stack code & Tries every change in the simulator first, then on the car \\
Maintainer & Our team, later another team & Builds the images, keeps the forks in sync, records design decisions \\
\bottomrule
\end{tabularx}

\subsubsection{Use Cases}
\begin{takeaway}
\textbf{UC-001 --- Map the room}\\
\textbf{Actor:} Operator\\
\textbf{Precondition:} car in the room, stack running, joystick connected\\
\textbf{Main flow:} 1.~Operator drives the car around the room by joystick; 2.~SLAM builds the room map; 3.~Operator saves the map\\
\textbf{Postcondition:} a saved map of the room (not of the track)
\end{takeaway}

\begin{takeaway}
\textbf{UC-002 --- Drive laps on the track centre seen by the camera}\\
\textbf{Actor:} Operator\\
\textbf{Precondition:} car on the track, stack running\\
\textbf{Main flow:} 1.~Operator holds the deadman; 2.~line detection publishes the track centre as a path in the car frame; 3.~the path follower drives along it\\
\textbf{Postcondition:} full laps without leaving the lines; releasing the deadman stops the car
\end{takeaway}

\begin{takeaway}
\textbf{UC-003 --- Drive laps on a path in the map}\\
\textbf{Actor:} Operator\\
\textbf{Precondition:} saved room map, a saved path in map coordinates \tbd[OI-03]\\
\textbf{Main flow:} 1.~Operator gives the initial pose; 2.~localization tracks the car in the map; 3.~the path is published once; 4.~the path follower drives along it\\
\textbf{Postcondition:} full laps without leaving the lines
\end{takeaway}

\begin{takeaway}
\textbf{UC-004 --- Stop or take over}\\
\textbf{Actor:} Operator\\
\textbf{Precondition:} the autonomy is driving\\
\textbf{Main flow:} the Operator presses the emergency stop, releases the deadman, or drives with the joystick; or an obstacle appears in front of the car\\
\textbf{Postcondition:} the car stops or follows the joystick, whatever the autonomy sends
\end{takeaway}

\begin{takeaway}
\textbf{UC-005 --- Try it in the simulator first}\\
\textbf{Actor:} Developer\\
\textbf{Precondition:} the container on a Mac or PC\\
\textbf{Main flow:} 1.~Developer starts the simulator as the car; 2.~runs the same stack and launch as on the car; 3.~then runs it on the car unchanged\\
\textbf{Postcondition:} the behaviour seen in the simulator is the behaviour on the car, within the simulator's models
\end{takeaway}

\subsubsection{Detailed Functional Requirements}
\paragraph{Drivers (DRV)}
\begin{tabularx}{\textwidth}{Q}
\reqhead
FR-DRV-001 & The stack shall drive the motor and the servo through the VESC and publish wheel odometry from it. & V1 & S \\
FR-DRV-002 & The stack shall read the joystick and turn it into drive commands while the operator holds the enable button. & V1 & S \\
FR-DRV-003 & The stack shall publish the Pulsar SL450 scans and the Gemini 335L colour images, camera info and IMU through \code{orbbec\_camera}. & V1 & S \\
FR-DRV-004 & The stack shall publish the mounting transforms from \code{base\_link} to the LiDAR and to the camera, from measured values. & V1 & C \\
\bottomrule
\end{tabularx}

\paragraph{Mapping (MAP)}
\begin{tabularx}{\textwidth}{Q}
\reqhead
FR-MAP-001 & The stack shall build a map of the room with SLAM (for example \code{slam\_toolbox}) and save it. & V1 & S \\
FR-MAP-002 & The map shall be of the room, not of the track. & V1 & S \\
\bottomrule
\end{tabularx}

\paragraph{Localization (LOC)}
\begin{tabularx}{\textwidth}{Q}
\reqhead
FR-LOC-001 & The stack shall localize the car in the saved room map with AMCL, from odometry and LiDAR scans. Node choice \tbd[OI-07]. & V1 & S \\
FR-LOC-002 & The stack shall accept an initial pose from the operator. & V1 & S \\
FR-LOC-003 & The stack should fuse wheel odometry and the camera IMU into \code{odom} with an EKF. \tbd[OI-08] & V1 & S \\
\bottomrule
\end{tabularx}

\paragraph{Line detection (LIN)}
\begin{tabularx}{\textwidth}{Q}
\reqhead
FR-LIN-001 & The stack shall detect the tape lines in the camera image. Method \tbd[OI-06]. & V1 & S \\
FR-LIN-002 & The stack shall publish the track centre as a \code{nav\_msgs/Path}, recomputed from every camera frame, covering a short stretch ahead of the car. Frame \tbd[OI-02]. & V1 & S \\
\bottomrule
\end{tabularx}

\paragraph{Path following (PTH)}
\begin{tabularx}{\textwidth}{Q}
\reqhead
FR-PTH-001 & A PID controller shall drive the car on the lateral (and heading) error between the car's pose and a path. Form \tbd[OI-05]. & V1 & S \\
FR-PTH-002 & The path shall be a \code{nav\_msgs/Path} on one topic, in any frame: drawn by the user (map), from line detection (car frame), or the racing line of option A (map). & V1 & S \\
FR-PTH-003 & The pose shall be read from TF in the path's frame. A path in the car frame shall need no map pose. & V1 & S \\
FR-PTH-004 & The controller shall publish \code{/drive}. Any node with the same inputs (path, pose) and output shall be able to replace it, for example Pure Pursuit. & V1 & S \\
FR-PTH-005 & Offline paths (hand-drawn, racing line) shall be computed once, saved to a file and published once as a closed loop (latched). Open paths \tbd[OI-04]. & V1 & S \\
FR-PTH-006 & On a full loop the controller shall search the closest point only in a window after the last one it used; on a fragment it shall search all of it. & V1 & S \\
FR-PTH-007 & If an online path stops updating (the camera loses the tape), the controller shall stop the car. A latched offline path is never stale. & V1 & S \\
FR-PTH-008 & The stack shall provide a path drawn in the map. Source \tbd[OI-03]: teach and repeat suggested for V1, tape mapping later. & V1 & S \\
\bottomrule
\end{tabularx}

\paragraph{Safety (SAF)}
\begin{tabularx}{\textwidth}{Q}
\reqhead
FR-SAF-001 & A person with the joystick shall always be able to override the autonomy: the joystick has priority over \code{/drive}. & V1 & S \\
FR-SAF-002 & An emergency stop shall stop the car whatever is driving. Mechanism \tbd[OI-09]. & V1 & S \\
FR-SAF-003 & The autonomy shall drive only while the operator holds the deadman; releasing it shall stop the car. Mechanism \tbd[OI-09]. & V1 & S, C \\
FR-SAF-004 & Automatic emergency braking (AEB) shall compute the time to collision from the LiDAR and brake under a threshold, whatever is driving, while still letting a person drive the car away. & V1 & S \\
FR-SAF-005 & Each command input shall time out: a source that stops sending loses control of the car. & V1 & C \\
\bottomrule
\end{tabularx}

\paragraph{Simulator (SIM)}\label{sec:fr-sim}
\begin{tabularx}{\textwidth}{Q}
\reqhead
FR-SIM-001 & The simulator shall stand in for the car with the same topics, frames and launch, so the stack cannot tell which one it runs on. & V1 & S, C \\
FR-SIM-002 & Every safety feature (FR-SAF-001 to 004) shall work in the simulator as on the car. & V1 & S \\
FR-SIM-003 & The stack shall never read the simulator's ground truth; it is for debugging and scoring only. & V1 & S \\
\bottomrule
\end{tabularx}

\paragraph{Documentation (DOC)}
\begin{tabularx}{\textwidth}{Q}
\reqhead
FR-DOC-001 & The platform shall be documented: setup, how to run each part, and design decisions, so another team or the professor can use it without us. & V1 & S \\
\bottomrule
\end{tabularx}

\paragraph{Advanced topic (ADV)}
\begin{tabularx}{\textwidth}{Q}
\reqhead
FR-ADV-001 & \textbf{Option A, racing line.} The stack shall get the track boundaries in map coordinates by mapping the tape lines seen by the camera into the room map. & V2 & S \\
FR-ADV-002 & \textbf{Option A.} It shall optimize the racing line (minimum curvature, then minimum lap time) within the car's limits, with TUM's \code{global\_racetrajectory\_optimization} as the reference, and compute the speed profile along it. & V2 & S \\
FR-ADV-003 & \textbf{Option A.} It shall track the line at speed with a controller that copes with high speed (tuned Pure Pursuit, or MPC). & V2 & S \\
FR-ADV-011 & \textbf{Option B, hybrid localization.} The stack shall fuse wheel odometry, IMU and LiDAR localization in the room map with an EKF (for example \code{robot\_localization}). & V2 & S \\
FR-ADV-012 & \textbf{Option B.} The detected tape lines shall correct the pose relative to the track, where the LiDAR is weak. & V2 & S \\
FR-ADV-013 & \textbf{Option B.} When one sensor drops out or gives a wrong reading, the pose shall stay usable. & V2 & S \\
\bottomrule
\end{tabularx}

\paragraph{Off-track report (OFF)}
\begin{tabularx}{\textwidth}{Q}
\reqhead
FR-OFF-001 & Optionally, a node may compare the car's footprint at the stack's pose with the tape lines and report when it believes the car is off the track. It only reports, never stops the car. \tbd[OI-11] & \tbd & S \\
\bottomrule
\end{tabularx}

\subsection{Non-Functional Requirements}
\subsubsection{Portability and variability}
\begin{tabularx}{\textwidth}{Q}
\reqhead
NFR-VAR-001 & The stack shall not depend on the track layout: any closed loop drawn as two boundary lines. & V1 & S \\
NFR-VAR-002 & The stack shall not depend on the room, which is mapped on site. & V1 & S \\
NFR-VAR-003 & Line detection should not depend on the floor or the tape colour, once those are known \tbd. & V1 & S \\
NFR-POR-001 & The same stack, images and launch shall run on the simulator (Mac, PC) and on the car (Orin Nano). & V1 & S \\
\bottomrule
\end{tabularx}

\subsubsection{Performance}
\begin{tabularx}{\textwidth}{@{}L{2.4cm} L{2.8cm} L{2.4cm} Y@{}}
\toprule
\hd{ID} & \hd{Requirement} & \hd{Target} & \hd{Justification} \\
\midrule
NFR-PER-001 & Top speed on the track & \tbd & Not set by the scope. The sensor choice was made for 5--8\,m/s (H). \\
NFR-PER-002 & Scan rate & 40\,Hz & At 8\,m/s a 10\,Hz LiDAR leaves 80\,cm between scans, too much for scan matching (H). \\
NFR-PER-003 & Camera frame rate & $\geq$ 60\,fps & At 8\,m/s, 30\,fps means 27\,cm travelled per frame (H). \\
NFR-PER-004 & Timestamp error & $\leq$ 10\,ms & 10\,ms at 8\,m/s is 8\,cm on the floor (H). \\
NFR-PER-005 & Command timeout & 0.2\,s per mux input & Today's \code{ackermann\_mux} configuration (C). \\
\bottomrule
\end{tabularx}

\subsubsection{Safety}
\begin{itemize}
  \item \textbf{NFR-SAF-001:} a person can always stop or take over the car, on the car and in the simulator (FR-SAF-001 to 005, FR-SIM-002).
  \item \textbf{NFR-SAF-002:} each safety feature is tested on the car (AC-V1-04) before autonomous laps on the car.
\end{itemize}

\subsubsection{Maintainability}
\begin{itemize}
  \item \textbf{NFR-MNT-001:} the forks stay upstream plus our upstreamable fixes; our code lives in \code{polimi\_*} packages in our repo.
  \item \textbf{NFR-MNT-002:} pinned dependencies are bumped on purpose, never by following a branch.
\end{itemize}

\subsection{Summer school exercises (reference)}\label{sec:exercises}
The exercises go from the simplest to the most advanced; each builds on the ones before. The scope says each must run in the simulator first, then on the car. Whether that makes them requirements is open \tbd[OI-12]: the suggestion is reference only, with the simulator covering what the base platform's acceptance needs.

\begin{tabularx}{\textwidth}{@{}L{4.2cm} Y L{2.6cm}@{}}
\toprule
\hd{Exercise} & \hd{Goal} & \hd{Platform part} \\
\midrule
\endhead
1. Trajectory tracking control (controllo e inseguimento traiettoria) & Follow a given path with a feedback controller. 1a with absolute localization (map path), 1b with computer vision (camera path). Same controller, only the reference changes. & FR-PTH \\
2. Line recognition (riconoscimento delle linee) & Detect the tape lines in the image and give their position on the floor, relative to the car. Input of 1b. & FR-LIN \\
3. Localization on a map (AMCL) & Map the room with SLAM, then localize in it with a particle filter. Input of 1a. & FR-MAP, FR-LOC \\
4. Planning (pianificazione) & Compute the path and speed profile, introduction to the racing line. & FR-ADV (A) \\
5. Hybrid localization (localizzazione ibrida) & Fuse odometry, IMU, LiDAR and camera with an EKF. & FR-LOC-003, FR-ADV (B) \\
\bottomrule
\end{tabularx}

% ============================================================================
\section{System Design}
% ============================================================================
\subsection{Architecture Overview}
The design is in \code{TDD\_F1tenth}. In short: the stack has three layers. The logic (SLAM, localization, line detection, path following) and a common layer (command mux and safety) are the same on car and simulator. Below them, exactly one hardware layer serves a fixed set of ROS topics and TF frames: the car's drivers, or the simulator. The stack talks to the hardware only through that set.

\subsection{System Components}
\begin{tabularx}{\textwidth}{@{}L{3.5cm} Y@{}}
\toprule
\hd{Component} & \hd{Requirements} \\
\midrule
Drivers & FR-DRV \\
Mapping, localization & FR-MAP, FR-LOC \\
Line detection & FR-LIN \\
Path following & FR-PTH \\
Mux and safety & FR-SAF \\
Simulator & FR-SIM \\
Advanced topic & FR-ADV \\
\bottomrule
\end{tabularx}

% ============================================================================
\section{Implementation Guidelines}
% ============================================================================
\subsection{Technology Stack}
ROS~2 Jazzy, \code{rmw\_zenoh\_cpp}, Docker images per machine type (\code{arm} Apple Silicon, \code{x86} PC, \code{jet} Orin Nano), \code{f1tenth\_system} drivers, \code{orbbec\_camera}, \code{slam\_toolbox}, \code{nav2\_amcl}, \code{robot\_localization}. Details in \code{TDD\_F1tenth}.

\subsection{Development Standards}
Every change goes through a pull request; the shared branches are protected. Python through \code{uv}. Runbooks in \code{docs/software/}.

\subsection{Testing Strategy}
Every function runs in the simulator first, then on the car. The acceptance below is checked on both.

\subsubsection{Acceptance criteria}
\begin{tabularx}{\textwidth}{@{}L{2.2cm} Y L{2.6cm}@{}}
\toprule
\hd{ID} & \hd{Done when} & \hd{Covers} \\
\midrule
\endhead
AC-V1-01 & The car maps the room and localizes with AMCL. & FR-MAP, FR-LOC \\
AC-V1-02 & With the PID it drives full laps of the taped track without leaving the lines, on the track centre seen by the camera, on the car and in the simulator. & FR-LIN, FR-PTH, FR-SIM \\
AC-V1-03 & The same, on a path drawn in the map. & FR-PTH-008, FR-LOC \\
AC-V1-04 & Tested on the car: pressing the emergency stop, releasing the deadman, taking over with the joystick and an obstacle in front each stop the car or hand it to the person. & FR-SAF \\
AC-V2-A & Option A: the car laps faster on the optimized line than on the centre line without leaving the lines, measured as lap times and line crossings over several laps. & FR-ADV-001 to 003 \\
AC-V2-B & Option B: the fused pose is more accurate and more robust than LiDAR-only localization, measured against a ground truth (the simulator, or marks on the floor) at increasing speeds. & FR-ADV-011 to 013 \\
\bottomrule
\end{tabularx}

Leaving the lines is judged by the people watching: nothing in software measures it on the car, and there is no external tracking (C-005, FR-OFF-001).

% ============================================================================
\section{Project Timeline}
% ============================================================================
\begin{tabularx}{\textwidth}{@{}Y L{3cm} L{3cm}@{}}
\toprule
\hd{Phase} & \hd{Start} & \hd{End} \\
\midrule
Hardware selection & Sep 2026 & 7 Oct 2026 (BOM) \\
Simulator on Humble & Oct 2026 & done \\
Jazzy migration (OI-13) & \tbd & \tbd \\
Base platform (V1) & \tbd & \tbd \\
Advanced topic (V2) & after V1 runs & \tbd \\
Summer school & \tbd & \tbd \\
\bottomrule
\end{tabularx}

% ============================================================================
\section{Risks and Mitigation}
% ============================================================================
\begin{tabularx}{\textwidth}{@{}L{4cm} L{1.6cm} L{1.6cm} Y@{}}
\toprule
\hd{Risk} & \hd{Prob.} & \hd{Impact} & \hd{Mitigation} \\
\midrule
\endhead
Chassis is end of life, low stock & Medium & High & Several retailers in the BOM; Plan B parts list in the sourcing study (H) \\
Glare or floor reflections hide the tape & Medium & High & Early phone video of the real tape; method choice OI-06; global-shutter camera \\
Symmetric rectangular room confuses localization & Medium & Medium & Always give the initial pose (OI-07) \\
People or moved furniture change the room after mapping & High & Medium & Re-map on site (A-001); option B corrects from the tape \\
Camera pitch under braking breaks the floor projection & Medium & Medium & Floor plane from depth, or the IMU (H) \\
No USB passthrough on Docker Desktop for macOS: no joystick in the simulator & High & Medium & A stand-in for the joystick in the simulator \tbd \\
A gain from option A too small to measure in a small room & Medium & Medium & Decide OI-01 after V1 runs \\
\bottomrule
\end{tabularx}

% ============================================================================
\section{Open Items}
% ============================================================================
Decisions and checks the requirements leave open. \code{docs/scope.md} holds the options and suggestions until this document reaches v1.00.

\begin{tabularx}{\textwidth}{@{}L{1.4cm} Y L{2.6cm}@{}}
\toprule
\hd{ID} & \hd{Item} & \hd{Affects} \\
\midrule
\endhead
OI-01 & Advanced topic: option A or B. Held until the base platform runs. & FR-ADV \\
OI-02 & Frame of the line-detection path: car frame or \code{odom}. & FR-LIN-002 \\
OI-03 & Where the path drawn in the map comes from (teach and repeat, then tape mapping). & FR-PTH-008 \\
OI-04 & Open paths: stop at the end of the path. & FR-PTH-005 \\
OI-05 & Form of the PID (lateral and heading error, look-ahead). & FR-PTH-001 \\
OI-06 & Line detection method. & FR-LIN-001 \\
OI-07 & Localization node (\code{nav2\_amcl}, \code{slam\_toolbox}, \code{particle\_filter}). & FR-LOC-001 \\
OI-08 & EKF for \code{odom} in the base platform. & FR-LOC-003 \\
OI-09 & Safety with \code{ackermann\_mux} locks [VERIFY]. & FR-SAF-002, 003 \\
OI-11 & Off-track report (optional). & FR-OFF-001 \\
OI-12 & Exercises in the simulator: reference or required. & \cref{sec:exercises} \\
\bottomrule
\end{tabularx}

Decided: OI-10 (one image rebuild with \code{nav2\_amcl}, \code{robot\_localization}, \code{image\_transport\_plugins}, \code{cv\_bridge}; 22 Oct 2026) and OI-13 (Jazzy, C-020; 21 Oct 2026).

% ============================================================================
\section{Glossary and References}
% ============================================================================
\subsection{Definitions}
\begin{description}
  \item[30L] 30 e lode, the top grade.
  \item[AEB] Automatic emergency braking.
  \item[AMCL] Adaptive Monte Carlo localization, a particle filter on a known map.
  \item[Deadman] A button the operator holds for the autonomy to drive; releasing it stops the car.
  \item[EKF] Extended Kalman filter.
  \item[Mux] The command multiplexer (\code{ackermann\_mux}): forwards the highest-priority input still sending.
  \item[SLAM] Simultaneous localization and mapping.
  \item[VESC] The motor controller; it also drives the steering servo and gives wheel odometry.
\end{description}

\subsection{References}
\begin{enumerate}
  \item \code{docs/scope.md}, Project scope, 21 Oct 2026.
  \item \code{docs/hardware/LiDAR\_comparison\_F1TENTH.md}, 23 Sep 2026.
  \item \code{docs/hardware/Camera\_comparison\_F1TENTH.md}, 7 Oct 2026.
  \item \code{docs/hardware/Fiesta\_ST\_Rally\_VXL\_sourcing\_F1TENTH.md}, 24 Sep 2026.
  \item \code{hardware/bom/master\_bom.xlsx}, updated 7 Oct 2026.
  \item \code{docs/software/car.md}, car bringup.
  \item \code{TDD\_F1tenth}, the design that refines these requirements.
\end{enumerate}

\end{document}
